\item Un átomo aislado de cierto elemento emite luz de longitud de onda $\lambda_{m1}$ cuando el átomo cae de su estado con número cuántico $m$ a su estado fundamental de número cuántico $1$. El átomo emite un fotón de longitud de onda $\lambda_{n1}$ cuando cae de su estado excitado con número cuántico $n$ a su estado fundamental.
\begin{enumerate}
    \item Encuentre la longitud de onda de la luz radiada cuando el átomo hace una transición de su estado $m$ a su estado $n$.
    \item Demuestre que $k_{mn} = |k_{m1} - k_{n1}|$, donde $k_{ij} = 2\pi/\lambda_{ij}$ es el número de onda del fotón. Este problema ejemplifica el principio de combinación de Ritz, una regla empírica formulada en 1908.
\end{enumerate}